\subsection{Evolução das notas de corte por sistema de concorrência}

Para avaliar como a exigência mínima de ingresso evoluiu em cada sistema, definiu-se a nota de corte de um grupo como o menor Argumento Final entre os candidatos convocados na primeira chamada do primeiro semestre nas categorias daquele grupo. Os quatro grupos seguem a classificação adotada nas seções anteriores: Sistema Universal, Cotas para Negros, EP PPI e EP Não-PPI, estes dois últimos reunindo as faixas de renda R1 e R2. As categorias PcD foram excluídas do cálculo, pois dispõem de uma única vaga, frequentemente não preenchida, e o mínimo delas não representa o grupo. No subprograma 2018--2020, em que o ingresso foi unificado, considerou-se a chamada única. Cabe ressaltar que o Argumento Final é padronizado dentro de cada coorte (Equações 1 a 3); assim, a comparação entre subprogramas indica a posição relativa do último aprovado em relação aos demais candidatos do mesmo ano, e não um nível absoluto de conhecimento.

\begin{figure*}[htbp]
\centering
\includegraphics[width=\textwidth]{F4_notas_corte.png}
\caption{Nota de corte (menor Argumento Final entre os convocados da primeira chamada do primeiro semestre) por sistema de concorrência nos cursos de Computação. Faixas sombreadas: vigência do SISUnB e provas aplicadas durante a pandemia. Linha tracejada: Lei nº 14.723/2023.}
\label{fig:f4_cortes}
\end{figure*}

A Figura~\ref{fig:f4_cortes} evidencia comportamentos distintos entre os cursos. Em Ciência da Computação, a nota de corte do Sistema Universal passou de 93 a 104 pontos entre 2015--2017 e 2017--2019 para 136 a 144 pontos entre 2020--2022 e 2023--2025, uma elevação de cerca de 40 pontos. Esse movimento acompanha o aumento da procura pelo curso, que passou de 138 candidatos não eliminados em 2015--2017 para até 423 em 2020--2022, mantidas as mesmas 40 vagas anuais, e converge com a tendência de crescimento da seletividade a partir de 2021 apontada na Seção~V-D. Em Engenharia de Computação, por sua vez, o corte do Sistema Universal permaneceu estável, em torno de 120 pontos, desde 2020--2022.

Nos dois bacharelados, as notas de corte das cotas de Escola Pública situam-se, na maior parte da série, entre $-40$ e $+20$ pontos, ou seja, próximas da média da coorte e cerca de 100 a 140 pontos abaixo do Sistema Universal. A cota para Negros ocupa posição intermediária: em Ciência da Computação, seu corte subiu de 41 a 50 pontos entre 2015--2017 e 2018--2020 para 83 a 137 pontos a partir de 2019--2021, aproximando-se do Sistema Universal. Em Engenharia de Computação, o corte dessa cota oscila amplamente (de 21 a 143 pontos), reflexo da existência de apenas uma vaga por semestre. Esses resultados são compatíveis com a redução do $d$ de Cohen observada na Seção~V-B: à medida que os cotistas de alto rendimento se aproximam do Sistema Universal no desempenho médio, a exigência mínima de algumas cotas também se eleva.

Quanto aos fatores exógenos, todas as notas de corte dos bacharelados caem em 2018--2020 (por exemplo, 78 e 76 pontos no Sistema Universal de Ciência da Computação e de Engenharia de Computação, respectivamente). Como nesse subprograma as 40 vagas anuais foram preenchidas em uma única seleção, essa queda está associada sobretudo à unificação do ingresso, e não necessariamente a um pior desempenho no período pandêmico. Já a partir da Lei nº 14.723/2023, que passou a preencher primeiro as vagas de ampla concorrência, espera-se que os cotistas de maior nota migrem para o Sistema Universal e que o corte das cotas diminua. Esse efeito aparece em Engenharia de Computação, onde o corte de EP PPI atinge $-47$ e $-67$ pontos em 2022--2024 e 2023--2025, os menores valores da série. Em Ciência da Computação, contudo, não há padrão claro, e com apenas dois subprogramas e categorias de uma a duas vagas não é possível afirmar uma mudança de tendência. Por fim, na Licenciatura em Computação, a procura é próxima ou inferior à oferta de 46 vagas anuais (de 27 a 72 candidatos não eliminados) e quase todos os candidatos são aprovados, de modo que o corte corresponde, em geral, à nota do último candidato disponível e não expressa concorrência efetiva.

\subsection{Recuo da nota de corte nas chamadas suplementares}

A nota de corte da primeira chamada não coincide com a nota do último candidato efetivamente convocado, já que desistências e vagas não ocupadas geram chamadas sucessivas e a seleção para o segundo semestre. Para mensurar esse recuo, calculou-se a diferença entre a nota de corte da primeira chamada do primeiro semestre e o menor Argumento Final entre todos os convocados do grupo, considerando todas as chamadas e ambos os semestres (Figura~\ref{fig:f5_recuo}). Valores positivos indicam que a exigência mínima cedeu ao longo do processo.

\begin{figure*}[htbp]
\centering
\includegraphics[width=\textwidth]{F5_gap_nota_corte.png}
\caption{Recuo da nota de corte: diferença entre o corte da primeira chamada do primeiro semestre e o Argumento Final do último candidato convocado (todas as chamadas e ambos os semestres), por sistema de concorrência. Linha tracejada: Lei nº 14.723/2023.}
\label{fig:f5_recuo}
\end{figure*}

Nos dois bacharelados, o recuo está presente em todos os grupos, com medianas entre 20 e 82 pontos ao longo da série. Em Ciência da Computação, as medianas foram de 34 pontos no Sistema Universal, 49 na cota para Negros, 40 em EP PPI e 25 em EP Não-PPI; em Engenharia de Computação, de 38, 82, 48 e 20 pontos, respectivamente. A cota para Negros apresenta os maiores recuos (com picos de 129 pontos em Ciência da Computação, em 2021--2023, e de 119 pontos em Engenharia de Computação, em 2017--2019), pois, com uma única vaga por semestre, a desistência de um convocado transfere a vaga a um candidato posicionado bem abaixo na lista. No Sistema Universal, que conta com uma lista de espera extensa, o recuo é menor e mais estável, entre 29 e 49 pontos a partir de 2019--2021. Os valores mais elevados de Engenharia de Computação em 2015--2017 e 2016--2018 (122 e 162 pontos) coincidem com a vigência do SISUnB, quando a possibilidade de troca de curso após o resultado ampliava as desistências e o número de chamadas.

Na Licenciatura em Computação, o recuo nas cotas é praticamente nulo, pois todos os candidatos cotistas já são convocados na primeira chamada e não há lista de espera. Em contrapartida, o recuo no Sistema Universal é elevado (mediana de 68 pontos), o que decorre da migração das vagas reservadas não preenchidas para a ampla concorrência, ocupadas por candidatos de menor Argumento Final. Após a Lei nº 14.723/2023, não se observa alteração sistemática no recuo em nenhum dos cursos.

\subsection{Representação dos cotistas entre candidatos e aprovados}

Por fim, para avaliar a efetividade do acesso via cotas, a Figura~\ref{fig:f6_representacao} compara, para cada curso e subprograma, a proporção de cotistas entre os candidatos não eliminados no PAS~3, a proporção de cotistas entre os aprovados (convocados em alguma chamada) e a proporção de aprovados que ingressaram por vaga reservada. Considerou-se cotista todo candidato inscrito na cota para Negros ou em qualquer categoria de Escola Pública, incluindo PcD. Como referência, a figura também apresenta a fração das vagas de cada curso reservada às cotas, de 55\% nos bacharelados e 56,5\% na licenciatura.

\begin{figure*}[htbp]
\centering
\includegraphics[width=\textwidth]{F6_representacao_cotistas.png}
\caption{Proporção de cotistas entre os candidatos, proporção de cotistas entre os aprovados e proporção de aprovados por vaga de cota nos cursos de Computação. A linha indica a fração das vagas reservadas às cotas. Linha tracejada: Lei nº 14.723/2023.}
\label{fig:f6_representacao}
\end{figure*}

Nos bacharelados, os cotistas são minoria entre os candidatos, mas representam cerca de metade dos aprovados. Em Ciência da Computação, eles correspondem a entre 22\% e 38\% dos candidatos e a entre 47\% e 57\% dos aprovados, próximo dos 55\% de vagas reservadas, de modo que a reserva praticamente dobra sua presença entre os ingressantes (razão entre as proporções de 1,4 a 2,4). Em Engenharia de Computação, os cotistas aprovados ficaram abaixo da reserva até 2018--2020 (entre 32\% e 46\%), por insuficiência de candidatos cotistas, e passaram a 45\%--58\% a partir de 2019--2021. Considerando a nota de corte do Sistema Universal na Figura~\ref{fig:f4_cortes}, muito acima dos cortes das cotas, esses resultados indicam que, sem a reserva de vagas, a participação desses estudantes nos dois bacharelados seria substancialmente menor.

Em contrapartida, a proporção de cotistas entre os candidatos de Ciência da Computação recuou de 38\% em 2017--2019 para 22\% em 2023--2025, período em que a procura total pelo curso cresceu. Esse dado converge com a Seção~V-F: o aumento da demanda pela Computação concentra-se em candidatos do Sistema Universal, e a política de cotas mantém a representação entre os aprovados, mas não se reflete em maior participação dos cotistas na demanda inicial. Na Licenciatura em Computação, quase todos os candidatos são aprovados, e a proporção de cotistas entre os aprovados reproduz a dos candidatos (entre 15\% e 44\%), permanecendo sempre abaixo da reserva de 56,5\%; as vagas reservadas sobram por falta de candidatos e migram para a ampla concorrência. Nesse curso, portanto, a limitação está na demanda, e não no acesso.

Por fim, a Lei nº 14.723/2023 alterou a via de entrada, mas não a proporção total de cotistas aprovados. A diferença entre os cotistas aprovados e os aprovados por vaga reservada, que corresponde aos cotistas que ingressaram pelo Sistema Universal, permaneceu pequena nos bacharelados antes e depois da lei: no máximo 2 pontos percentuais em Ciência da Computação e até 5 pontos em Engenharia de Computação. Assim, poucos cotistas alcançam a nota de corte da ampla concorrência nesses cursos, e a nova regra modifica sobretudo o tipo de vaga ocupada, sem alterar de forma perceptível o número de cotistas que ingressam.
